\documentclass[../ppgc-tese.tex]{subfiles}
%\graphicspath{{\subfix{../images/}}}
\begin{document}

% =============================================================================
% Chapter source: "Magnitude, Not Meaning: What Activation-Space Adversarial
% Detection Measures in Vision Transformers" (interp_adversarial_image_paper,
% main.tex + appendix.tex, final version 2026-07-28).
%
% Narrative contract (Breno, 2026-09-10): state the hypothesis first, then
% present each experiment as the elimination of one objection to it. The
% structure must NOT become rhetoric: objections that are only refined, or that
% remain open, are labelled as such in the ledger (Table~\ref{tab:adv-ledger})
% and in the prose.
%
% Evidence-status legend, kept in comments next to each claim:
%   [EST]  established by data on disk in the research repo
%   [PRE]  pre-registered frozen gate / anti-condition (never relaxed)
%   [EXPL] exploratory run whose expectations were written down before execution
%   [EXT]  external cited claim
% Every number below is copied from main.tex/appendix.tex, whose \src trailers
% name the artifact; the trailer is preserved here as a "% src:" comment.
% =============================================================================

Sparse autoencoders (SAEs) promise interpretable features in vision models, and
a natural application is adversarial-attack detection: reading, from a model's
own activations, whether an input has been manipulated. This chapter asks not
whether such a detector can be built, since it plainly can, but what it
measures. The answer, for a Vision Transformer classifier probed with a linear
detector under a leakage-free protocol, turns out to be physical rather than
semantic. Adversarial optimization buys activation displacement cheaply: at
equal input norm, attacks arrive at the probed layer with $3$--$10\times$ the
displacement of benign corruptions, and even the minimal-norm attack that flips
every label lands at attack scale (Figure~\ref{fig:adv-amplification}). Once
benign perturbations are calibrated to move the activation as far as the attack
does, they become at least as detectable as the attack, at every layer and
across five seeds; what still separates perturbation families is low-level
texture, not adversarialness, and the SAE's interpretable features are not
load-bearing at any step. % [EST — preview of sec:adv-specificity and sec:adv-residue]
An attack built instead to minimize displacement reaches benign scale only at a
$34\times$ input-norm premium, on visibly corrupted images, and is still
detected; whether the probe there reads adversarial structure or the wake of
that outsized input perturbation is the one question this chapter leaves
open. % [EST — sec:adv-lowdx; open question owned there]

\begin{figure}[ht]
  \centering
  \includegraphics[width=0.9\linewidth]{figures/interp_adv/fig_amplification.pdf}
  \caption{Input norm versus what the detector reads. At the same input $L_2$,
  attacks arrive at \texttt{mlp\_10} of ViT-B/16 with $3$--$10\times$ the
  activation displacement $\lVert\Delta x\rVert$ of benign perturbations; the
  minimal-input-norm Carlini--Wagner attack flips every label from the smallest
  input budget yet still lands at attack-scale displacement; the steganographic
  control barely moves the representation at all. This figure is the preview
  of the chapter's turning point (Section~\ref{sec:adv-specificity}).}
  % src: figs/src/fig_amplification.py header lists the seven source artifacts
  \label{fig:adv-amplification}
\end{figure}

The chapter is organized as a defence. Section~\ref{sec:adv-hypotheses} states
the hypothesis under test, the rival hypothesis it was pre-registered against,
and the ledger of objections that the remaining sections address one by one.
Section~\ref{sec:adv-background} gives the background specific to this study.
Section~\ref{sec:adv-instrument} fixes the instrument. Section~\ref{sec:adv-gates}
reports the pre-registered test, and Sections~\ref{sec:adv-regime}
to~\ref{sec:adv-lowdx} work through the objections, ending with the one that
stays open. Section~\ref{sec:adv-reconciliation} reconciles the result with the
deployment literature, and Section~\ref{sec:adv-closing} states the
limitations and the reading for the thesis.

% =============================================================================
\section{Two hypotheses and a confound}
\label{sec:adv-hypotheses}

Adversarial examples remain among the most studied failure modes of deep vision
models: small input perturbations that flip a confident prediction
\cite{goodfellow2015fgsm,costa2024adversarial_survey}. A recurring hope is that
the model's own internal representations betray the attack, that somewhere in
activation space there is a signal for ``this input is adversarial'', and that
SAEs, by decomposing activations into interpretable features, can expose it
\cite{rao2024discover_then_name,pach2025sae_monosemantic_vlm}. On this view,
adversarial detection becomes a form of introspection: the model, read through
its features, ``knows'' it is under attack. % [EXT/context]

This hope rests on a confound that is rarely isolated. Every adversarial example
is two things at once: it is \emph{adversarial}, crafted to change the output,
and it is a \emph{perturbation}, so the activation has moved off its clean
location. A detector that fires on adversarial inputs may be reading either
property, the adversarial structure or merely the displacement. The two are
entangled in any standard evaluation, because the only negatives are clean
inputs, which differ from adversarial ones in both respects at once.
Disentangling them requires a control that standard detection setups omit: a
benign perturbation that moves the representation as far as the attack does.

\subsection{The hypothesis defended and the hypothesis pre-registered}
\label{sec:adv-two-hypotheses}

Two hypotheses compete for the same observations, and the chapter is honest
about which one was written down first.

\paragraph*{The introspection hypothesis ($H_{M}$).}
Sparse-feature signatures of an attacked input carry a signal that is specific
to the adversarial class and invariant to the perturbation's magnitude. This is
the hypothesis the thesis proposal pre-registered, operationalized as two frozen
gates: \emph{H1 (magnitude invariance)}, detection AUC should not track the
attack budget $\varepsilon$; and \emph{H2 (adversarial specificity)}, the attack
should be markedly more detectable than a benign perturbation of matched
magnitude. The gate definitions are given in Section~\ref{sec:adv-gates} and
were fixed before any run. % [PRE]

\paragraph*{The displacement hypothesis ($H_{D}$).}
To a linear probe on mean-pooled activations, the detectability of a
perturbation is a function of the activation-space displacement
$\lVert\Delta x\rVert$ that the perturbation induces at the probed layer.
Adversarialness has no privileged role beyond being one cause of that
displacement. The instrument scope, linear and mean-pooled, is part of the
statement rather than a caveat added later: $H_{D}$ makes no claim about what a
richer probe would see (Section~\ref{sec:adv-instrument-scope}).

The pre-registered gates are the contest between the two. If $H_{M}$ holds, the
gates pass; if $H_{D}$ holds, they fail in a specific way, with AUC rising
monotonically with budget and benign controls of equal displacement detected as
readily as the attack. Section~\ref{sec:adv-gates} reports that they fail. The
failure of $H_{M}$ is, however, only weak evidence \emph{for} $H_{D}$: the
frozen-budget runs are saturated, and a rival reading survives every gate
failure. The rest of the chapter is therefore organized around $H_{D}$ and the
objections that can be raised against it, each answered by an experiment
designed for that purpose.

\subsection{The specificity ladder}
\label{sec:adv-ladder}

It helps to organize ``how specific is the detected signal?'' as a ladder with
three degrees (Figure~\ref{fig:adv-ladder}). At \emph{Degree~1} the signal is
specific to the \emph{attack method}: the features that expose one attack
family differ from those exposing another. At \emph{Degree~2} it is specific to
the \emph{adversarial class}: adversarial features differ from those of a benign
perturbation of equal magnitude. At \emph{Degree~3} it is \emph{generic
displacement}: the detector fires for anything that moves the activation off
its clean location, attack or not. A magnitude-matched benign control is
precisely the instrument that separates Degree~2 from Degree~3, and the chapter
will show that even that control is not enough unless the match is made on the
axis the detector reads. $H_{M}$ lives at Degree~2 or above; $H_{D}$ is the
claim that the detector sits at Degree~3.

\begin{figure}[ht]
  \centering
  \includegraphics[width=0.85\linewidth]{figures/interp_adv/fig_ladder.pdf}
  \caption{The specificity ladder. Each control descends one level: the
  cross-attack analysis of SAEgis \cite{wang2026saegis} leaves Degree~1; a
  matched input-$L_2$ control appears to stop at Degree~2 with a residual
  $\dauc$, but that difference prices in amplification; the
  matched-$\lVert\Delta x\rVert$ control completes the descent to Degree~3, with
  texture as the remaining, non-adversarial, family axis.}
  \label{fig:adv-ladder}
\end{figure}

\subsection{The ledger of objections}
\label{sec:adv-ledger}

Table~\ref{tab:adv-ledger} lists every objection to $H_{D}$ that the study
confronted, the experiment built to answer it, and the verdict. Three entries
deserve attention precisely because they do not read as clean eliminations.
Objection~1 was answered by an experiment that \emph{strengthened} the
objection, producing the strongest evidence against $H_{D}$ in the study and
forcing the axis change that organizes everything downstream. Objection~5 was
not eliminated but forced a refinement of $H_{D}$: what remains at equal
displacement is texture. Objection~7 remains open. The ``origin'' column
separates the pre-registered gates from the exploratory experiments whose
expectations were written down before execution but after the gates had failed;
a committee should read the latter as designed follow-ups, not as
pre-registered tests.

\begin{table}[p]
  \centering
  \caption{The ledger of objections to the displacement hypothesis $H_{D}$.
  Status: \emph{eliminated}, \emph{refined} (the hypothesis had to be amended),
  or \emph{open}. Origin: PRE, pre-registered frozen gate; EXPL, exploratory
  experiment with expectations pre-stated; EXT, external evidence.}
  \label{tab:adv-ledger}
  \footnotesize
  \setlength{\tabcolsep}{3pt}
  \renewcommand{\arraystretch}{1.15}
  \begin{tabular}{@{}c >{\raggedright\arraybackslash}p{0.25\linewidth} >{\raggedright\arraybackslash}p{0.24\linewidth} >{\raggedright\arraybackslash}p{0.19\linewidth} >{\raggedright\arraybackslash}p{0.17\linewidth} >{\raggedright\arraybackslash}p{0.06\linewidth}@{}}
    \toprule
    \# & Objection to $H_{D}$ & Experiment & Verdict & Status & Orig. \\
    \midrule
    0 & $H_{M}$ is true: the signal is magnitude-invariant and adversarial-specific & H1 $\varepsilon$-sweep and H2 matched-$L_2$ benigns at the frozen budget, FGSM and PGD (\S\ref{sec:adv-gates}) & Both gates fail on every criterion; gap measured against a ceiling & $H_{M}$ eliminated in the tested regime; saturation weakens the test & PRE \\
    1 & The gates only fail because AUC is saturated & Decomposition and H2 re-run at $\varepsilon{=}0.004$ with the $L_2$ match verified (\S\ref{sec:adv-regime}) & A residual $\dauc{=}0.199$ appears and replicates & \textbf{Not eliminated}: strengthens the objection and motivates the axis change & EXPL \\
    2 & AUC rises with $\varepsilon$ because adversarial structure grows with $\varepsilon$ & Success-conditioned AUC (\S\ref{sec:adv-regime}) & Successful and failed attacks equally detectable & Eliminated & EXPL \\
    3 & The residual $\dauc$ is Degree-2 specificity & Per-example $\lVert\Delta x\rVert$, then benigns matched on $\lVert\Delta x\rVert$, five seeds, all layers (\S\ref{sec:adv-specificity}) & $\dauc\leq0$ at every layer & Eliminated; the hinge of the chapter & EXPL \\
    4 & The SAE is load-bearing: the gap lives in $z$, not in $a$ & Residual decomposition, signal zoo, logit-scale anti-condition (\S\ref{sec:adv-specificity}) & $a$ detects better than $z$; the gap lives on the manifold and vanishes at matched $\lVert\Delta x\rVert$ & Eliminated for detection & EXPL, PRE \\
    5 & If only magnitude mattered, all families would collapse onto one score curve & Collapse test with nested $F$-test (\S\ref{sec:adv-residue}) & Collapse fails, but the separating axis is texture & \textbf{Refined}: magnitude plus texture & EXPL \\
    6 & The probe reads the direction that flips the label & Projection ablation onto the head subspace (\S\ref{sec:adv-residue}) & Most energy and all of the signal in $\Delta x_{\perp}$ & Eliminated, with an enrichment caveat & EXPL \\
    7 & Adversarial and high-$\lVert\Delta x\rVert$ are confounded by construction & Minimal-input-norm CW, then a purpose-built minimal-displacement attack (\S\ref{sec:adv-lowdx}) & Cell occupiable at a $34\times$ input premium; still detected, unlike JPEG at equal $\lVert\Delta x\rVert$ & \textbf{Open}: asymmetry reported, not interpreted & EXPL \\
    8 & Stego at chance shows a blind probe, not a small displacement & Token-level probe and measured $\lVert\Delta x\rVert$ (\S\ref{sec:adv-regime}) & Watermark tokens light up; pooled displacement is $1\%$ of the norm & Eliminated; the null is instrument-scoped & EXPL \\
    9 & Then any off-manifold distance should reproduce the probe & Mahalanobis on the same activations (\S\ref{sec:adv-residue}) & $0.710$ against $0.989$ & \textbf{Refined}: a learned direction, not a generic distance & EXPL \\
    10 & The numbers are single-seed or pipeline artifacts & Clean-vs-clean and permutation nulls, paired bootstrap, five-seed replication, two documented bugs (\S\ref{sec:adv-protocol}) & Nulls at $0.500$; hinge replicates & Eliminated for the hinge; remaining single-seed scope declared & EXPL \\
    11 & SAEgis shows that SAE-based detection works & Reconciliation and the ladder (\S\ref{sec:adv-reconciliation}) & Three corroborations; only the strong reading is contested & Repositioned & EXT \\
    \bottomrule
  \end{tabular}
\end{table}

\subsection{Relation to the thesis question}
The study answers the diagnostics research question of the thesis:
under pre-registered tests, when do sparse-feature signatures isolate a
structurally specific perturbation and when do they collapse into change
detection? The chapter's answer is that, for adversarial perturbations that
optimize misclassification alone, they collapse into change detection, and the
change is measurable as a single physical variable. The methodological
contribution is the control that made that answer visible: matching benign
perturbations on the displacement the detector reads, not on the input.

% =============================================================================
\section{Background specific to this study}
\label{sec:adv-background}
% Condensed from the paper's Related Work (four paragraphs). SAE fundamentals
% are in the Background chapter; only what that chapter lacks is kept here.
% All paraphrase; no verbatim sentences from cited work.

\paragraph*{Detecting adversarial examples in representation space.}
Small crafted perturbations that flip a model's prediction
\cite{goodfellow2015fgsm,madry2018pgd,carlini2017cw_attack} prompted a
recurring proposal: read the attack off internal statistics, via auxiliary
classifiers \cite{metzen2017detecting_adversarial}, density and uncertainty
estimates \cite{feinman2017adversarial_artifacts}, input squeezing
\cite{xu2018feature_squeezing}, local intrinsic dimensionality
\cite{ma2018lid_adversarial}, reconstruction defences \cite{meng2017magnet},
or class-conditional Mahalanobis distances \cite{lee2018mahalanobis}, the
baseline run in Section~\ref{sec:adv-residue}. A counter-literature showed most
of these fail against adaptive adversaries
\cite{carlini2017bypassing_detection,sharma2018bypassing_feature_squeezing,athalye2018obfuscated_gradients,tramer2020adaptive_attacks},
detection being provably about as hard as robust classification
\cite{tramer2022detection_hardness}, and rigorous-evaluation guidelines followed
\cite{carlini2019evaluating_robustness}, which the pre-registered gates of this
chapter inherit. ViT-specific robustness work, from CNN comparisons to attention
patch attacks and architectural defences
\cite{bhojanapalli2021vit_robustness,naseer2021vit_properties,lovisotto2022attention_harmful,huang2023patchcensor},
acts at the architecture level, leaving open the representation-level seam this
detector occupies. What this literature rarely asks is what a working detector
\emph{reads}: while some detectors scale a noise baseline to the attack's
magnitude \cite{feinman2017adversarial_artifacts,ma2018lid_adversarial}, the
benign negatives are not verified to match each attack family, let alone matched
in activation displacement, so displacement and adversarialness stay
confounded. That verified per-family match, and above all the
displacement-matched regime, is this chapter's subject. % [EXT]

\paragraph*{Sparse autoencoders for vision models.}
The Background chapter (Chapter~\ref{chpt:background}) introduces SAEs as
overcomplete dictionaries of sparse, often monosemantic features. On vision
backbones the extraction machinery has matured rapidly: concept discovery and
naming
\cite{rao2024discover_then_name,gorton2024curve_detectors_sae,pach2025sae_monosemantic_vlm,lim2025sae_visual_concepts,zaigrajew2025clip_hierarchical_sae,stevens2025testable_vision_sae},
dictionary stability and cross-model alignment
\cite{fel2025archetypal_sae,thasarathan2025universal_sae}, and causal steering
of ViT features
\cite{joseph2025steering_clip_sae,han2025causal_sae_vision}. This chapter takes
that machinery as given and uses it as a measuring instrument: the clean-trained
SAE defines the on-/off-manifold split of Equation~\eqref{eq:adv-decomp}, and
the question is whether its features carry attack-specific information. % [EXT]

\paragraph*{SAEs meet adversarial inputs.}
The intersection is nascent. SAEgis \cite{wang2026saegis} is, to our knowledge,
the only SAE-based adversarial \emph{detector}, a deployment-oriented firewall
for vision--language models; Section~\ref{sec:adv-reconciliation} reconciles
this chapter's characterization with its findings. Adjacent work removes unsafe
concepts with SAEs \cite{geng2025sae_unsafe_concepts}, shows that SAE concept
representations are themselves attackable \cite{li2026sae_concept_robustness},
the concern behind the adaptive-attack item left for future work, and links
superposition to adversarial vulnerability \cite{bereska2025superposition_lossy}.
None of these supplies the magnitude- or displacement-matched control
contributed here. % [EXT]

\paragraph*{Sensitivity, local linearity, and amplification.}
Whether a small input perturbation grows into a large internal displacement is a
question of network sensitivity, and it has a literature. The linear view of
adversarial examples already predicts that a norm-bounded input perturbation
induces an activation change that scales with dimensionality
\cite{goodfellow2015fgsm}; the layer-to-layer amplification is governed by the
Lipschitz constant, which can grow sharply with depth unless constrained
\cite{cisse2017parseval}; and perturbations imperceptible in pixel space are
known to surface as substantial noise in intermediate feature maps
\cite{xie2019feature_denoising}. A smaller line studies transformers through
Jacobians and local linearity: activation steering via local-linear models
\cite{skifstad2026local_linearity}, layer-wise Jacobian regularization
\cite{martnishn2026dreg_jacobian}, Fisher-spectral robustness measures
\cite{zhang2026fisher_robustness}, and Lipschitz bounds on attention
\cite{yudin2025attention_lipschitz}. That work seeks to bound, denoise, or
exploit the amplification; this chapter \emph{measures} it per example and asks
what a detector reading the resulting displacement responds to. The
head-subspace ablation of Section~\ref{sec:adv-residue} is the simplest exact
member of the Jacobian family, with $J{=}W$ at the head input. % [EXT]

% =============================================================================
\section{The instrument}
\label{sec:adv-instrument}

\subsection{Models and activations}
\label{sec:adv-models}
The system is a Vision Transformer \emph{classifier}, not a vision--language
model: an ImageNet-1k \cite{deng2009imagenet} pretrained ViT-B/16 backbone
\cite{dosovitskiy2021vit} with a linear-probe classification head trained on
CIFAR-10 \cite{krizhevsky2009cifar} for three epochs (validation accuracy
$0.952$). % src: config/model/vit_b16.yaml; DESIGN_CHOICES.md §D1 [EST]
ViT-B/16 is the primary model and ViT-L/16 the scale check (head validation
accuracy $0.9548$). % src: config/model/vit_l16.yaml [EST]
CIFAR-10 images are upsampled to $224\times224$ and normalized with ImageNet
statistics throughout, consistent with the pretrained backbone; all runs use
seed~42 unless a multi-seed protocol is stated.
% src: config/data/cifar10.yaml; DESIGN_CHOICES.md §D1/§D7; config/base.yaml [EST]

ViT-B/16 processes each image as $197$ tokens ($196$ patches plus a class token)
through $12$ blocks of width $768$; ViT-L/16 has $24$ blocks of width $1024$.
The probe reads the MLP-output activations of blocks $\{0,3,6,9,10,11\}$ (and
the final pair $\{22,23\}$ for ViT-L/16), with block~10 canonical throughout.
% src: config/base.yaml; experiments/2026-06-22_sae-residual-decomp/params.yaml [EST]
Following SAE convention, every activation is \emph{mean-pooled over the $197$
tokens}: a deliberate simplification that erases token-localized structure and
bounds what the probe can see. Section~\ref{sec:adv-instrument-scope} states
the bound, and Section~\ref{sec:adv-regime} measures it in the one case where it
matters. % [EST/PRE — instrument scope]

\subsection{Sparse autoencoders and the residual decomposition}
\label{sec:adv-sae}
The whole chapter rests on one decomposition of the activation,
\begin{equation}
  a \;\approx\; \hat{a} + r,
  \qquad \hat{a} = \mathrm{dec}(z), \quad z = \mathrm{ReLU}(\mathrm{enc}(a)),
  \qquad r = a - \hat{a},
  \label{eq:adv-decomp}
\end{equation}
where $a$ is the mean-pooled raw activation at the probed block (the $a_\ell$ of
Section~\ref{sec:adv-dx-axis}, with the block index dropped when it is fixed),
$z$ the sparse latents, $\hat{a}$ the reconstruction, and $r$ the residual. The SAE is a vanilla $L_1$ SAE (input
dimension $768$, $1024$ latents, $L_1$ penalty $0.1$; Adam, learning rate
$10^{-3}$, batch $256$, ten epochs). % src: config/sae/default.yaml [EST]
Crucially, it is trained on \emph{clean activations only}, so $\hat{a}$ projects
onto the clean-data manifold and $r$ is, by construction, the off-manifold part
the clean dictionary cannot explain. Detection AUC is evaluated for the four
vector representations $a$, $z$, $\hat{a}$, $r$ and three scalar controls
($\lVert r\rVert_2$, $\lVert\hat{a}\rVert_2$, $\lVert z\rVert_1$), plus a
sparsity-matched TopK SAE ($k{=}128$) as a parametrization check.
% src: experiments/2026-06-22_sae-residual-decomp/params.yaml; config/sae/topk.yaml; D3b_RESULTS.md §3.2 [EST]

The logic of the decomposition is a competition. If interpretable SAE features
carried the adversarial signal, the latents $z$ should win; if the signal is
mere off-manifold displacement, the residual $r$, or even its scalar norm,
should already capture most of it. Because the SAE is a lossy bottleneck, $z$
can only discard detection information relative to $a$, never manufacture it,
so ``$z$ beats $a$'' would be the signature of a feature-level signal and
``$a \geq z$'' the signature of the SAE not being load-bearing.

\paragraph*{SAE health gates.}
Every $z$, $\hat{a}$, or $r$ comparison is gated on SAE health: fraction of
variance explained (FVE) $\geq 0.9$, dead-feature fraction $\leq 30\%$, and
$L_0$ in a sane band. % [PRE — health gate]
The canonical-layer SAE is healthy (FVE $0.938$, $L_0\approx135$, dead fraction
$3$--$5\%$); away from it, health varies by run, and in practice only
\texttt{mlp\_10} and \texttt{mlp\_11} ever pass (early layers reach FVE
$0.078$--$0.352$ in the six-layer decomposition).
% src: magnitude_sweep_metrics.csv; h2_matched_dx_summary.json; 2026-07-21_decomp-6layer [EST]
% Internal note carried over from the paper: FVE at mlp_0 differs between runs
% (0.75 in the 2026-06-22 decomposition vs 0.08 in the 2026-07-21 six-layer run);
% unexplained fit difference, flagged rather than resolved.
Tables therefore print the per-run FVE, sub-gate cells are marked, and prose
claims about latent-dependent quantities are restricted to passing layers;
claims about the raw activation $a$ are exempt.

\subsection{The displacement axis and the design matrix}
\label{sec:adv-dx-axis}
The variable that $H_{D}$ refers to is not the input-space perturbation budget but the
\emph{activation-space displacement} at the probed layer. For source image $i$
at layer $\ell$, let $a_\ell(\cdot)$ be the mean-pooled activation and
\begin{equation}
  \Delta x_\ell(i) \;=\; a_\ell(\tilde{x}_i) - a_\ell(x_i),
  \label{eq:adv-dx}
\end{equation}
the displacement of the perturbed input $\tilde{x}_i$ from its own clean
activation; $\lVert\Delta x\rVert_2$ is measured \emph{per example, per layer},
in every run. Input-space $L_2$ is a knob turned when constructing
perturbations; $\lVert\Delta x\rVert$ is what a probe reading layer $\ell$ can
possibly respond to. The two must not be conflated, because the network's
input-to-activation amplification is strongly family-dependent
(Section~\ref{sec:adv-specificity}); the amplification is itself a reported
quantity, not an assumption. % [EST]

The design space is a matrix (Table~\ref{tab:adv-design}): functional
adversarialness (does the perturbation flip the label?) crossed with activation
displacement. Its cells are populated by FGSM, PGD, and CW (adversarial, high
$\lVert\Delta x\rVert$), the four benign families (non-adversarial, high or
matched), the steganographic control (non-adversarial, low; attack success rate
$\approx 0.0007$, so not functionally adversarial),
% src: experiments/2026-06-29_h2-v2-eps0004/outputs/h2_crossval/metrics.json (asr block) [EST]
and, crucially, a cell no standard attack occupies: a successful attack with
low displacement. $H_{D}$ predicts that detectability is a function of the
displacement column only; any perturbation that breaks row-invariance breaks
the hypothesis. Section~\ref{sec:adv-lowdx} attempts to occupy the empty cell
by construction.

\begin{table}[ht]
  \centering
  \caption{The design matrix: functional adversarialness $\times$ measured
  activation displacement. $H_{D}$ predicts detectability depends on the column
  only. CW was the pre-registered \emph{candidate} for the top-right cell
  (minimal input norm); measured, it lands in the top-left, because the network
  amplifies it (Section~\ref{sec:adv-specificity}).}
  \label{tab:adv-design}
  \begin{tabular}{lcc}
    \toprule
    & High $\lVert\Delta x\rVert$ & Low $\lVert\Delta x\rVert$ \\
    \midrule
    Flips label      & FGSM, PGD, CW & \emph{empty} (Section~\ref{sec:adv-lowdx}) \\
    Does not flip    & four benign families & stego (control) \\
    \bottomrule
  \end{tabular}
  % src: EXPERIMENT_PLAN_ATTACKS.md §0
\end{table}

\subsection{Attacks, controls, and two matching regimes}
\label{sec:adv-attacks}

\paragraph*{Attacks.}
The attacks are untargeted FGSM \cite{goodfellow2015fgsm} with an $L_\infty$
budget, swept over the eight budgets from $\varepsilon{=}0.001$ to
$\varepsilon{=}0.032$ listed in Table~\ref{tab:adv-h1};
% src: config/attack/fgsm.yaml; experiments/2026-06-23_corrected-rerun/params.yaml [EST]
PGD \cite{madry2018pgd} with $20$ steps over the same grid, pre-registered as
the most rigorous H1 test;
% src: config/attack/pgd.yaml; experiments/2026-07-02_pgd-h1h2/outputs/h1_sweep/magnitude_sweep_metrics.csv [PRE — run]
and Carlini--Wagner $L_2$ \cite{carlini2017cw_attack}, the minimal-input-norm
attack, as the candidate for the low-displacement adversarial cell (attack
success rate $1.000$ on all $1600$ source images).
% src: experiments/2026-07-21_cw-attack/outputs/full/asr_cw.json [EXPL — E7]
Realized per-image input $L_2$ is measured in every run and checked against the
$\varepsilon$ arithmetic: FGSM has $L_2=6.85$ at $\varepsilon{=}0.004$ and
$20.51$ at $\varepsilon{=}0.012$ in normalized space.
% src: 2026-06-29_h2-v2-eps0004 metrics.json (l2_report); RESEARCH_LOG.md 2026-04-23 [EST]
Attack success rates (ASR) accompany every detection number: FGSM $0.762$ at
$\varepsilon{=}0.004$; PGD $0.922$ at $\varepsilon{=}0.002$; CW $1.000$.
% src: same metrics.json; 2026-07-02_pgd-h2-eps0002 metrics.json; asr_cw.json [EST]

\paragraph*{The low-displacement control.}
Steganography, a faint text watermark at opacity $0.03$, carries a near-zero
input payload ($L_2=2.24$, $0.33\times$ FGSM at $\varepsilon{=}0.004$).
% src: 2026-06-29_h2-v2-eps0004 metrics.json (l2_report) [EST]
Its measured displacement is the smallest of any family
($\lVert\Delta x\rVert=0.09$ at \texttt{mlp\_10}, under $1\%$ of the clean
activation norm), so it anchors the low-displacement column of the design
matrix by measurement, not by assumption.
% src: experiments/2026-07-23_jpeg-matched-l2-dx/outputs/displacement_summary.json [EST]

\paragraph*{Benign families and two matching regimes.}
To separate ``adversarial'' from merely ``perturbed'', four benign perturbation
families are used (Gaussian noise, Gaussian blur, JPEG compression, and random
pixel dropout), calibrated to match the attack in one of two regimes. Under
\emph{matched input-$L_2$}, each family's realized per-image input $L_2$ falls
within $\pm20\%$ of the attack's, verified in-run (ratios $0.999$--$1.008$ of
FGSM's at $\varepsilon{=}0.004$, with clean-alignment guards); non-matched
variants at $L_2\approx17$--$21$ are reported but never gated.
% src: 2026-06-29_h2-v2-eps0004 metrics.json [EST — closes audit gap R-C1]
Under \emph{matched $\lVert\Delta x\rVert$}, a model-in-the-loop calibration
searches each family's parameter until its mean displacement at the canonical
layer matches FGSM's (target $1.994$ at \texttt{mlp\_10}; realized
$1.78$--$2.06$, all within tolerance), the regime the displacement axis demands;
Section~\ref{sec:adv-specificity} shows the two regimes give sharply different
answers. % src: 2026-07-22_h2-matched-dx/outputs/calib/dx_calibration.json [EST — E1]

\subsection{Leakage-free protocol and statistical controls}
\label{sec:adv-protocol}
Every detection AUC uses one shared routine: grouped $K$-fold cross-validation
($K{=}5$) keyed to the \emph{source image}, so all perturbations of image~$i$
(clean, FGSM, stego, benign) share a fold and no detector is tested on a
near-duplicate of a training example. Features are standardized on the training
fold only; the detector is a converged logistic regression (LBFGS, $\leq 1000$
iterations); each binary task pits clean against \emph{one} family, never a
pool, at $N{=}3200$ ($1600$ clean, $1600$ perturbed).
% src: src/adv_interp/detection/evaluation.py; 2026-06-29_h2-v2-eps0004 h2_crossval/metrics.json [EST]

Throughout, specificity is summarized by the signed quantity
\begin{equation}
  \dauc \;=\; \mathrm{AUC}(\text{clean vs.\ attack})
  \;-\; \max_{b}\;\mathrm{AUC}(\text{clean vs.\ benign}_{b}),
  \label{eq:adv-dauc}
\end{equation}
with the maximum taken over the magnitude-matched benign families $b$. It is
positive when the attack is more detectable than every benign perturbation of
equal magnitude, and negative when some benign family is more detectable than
the attack; nothing constrains its sign. The frozen H2 gate of
Section~\ref{sec:adv-gates} is defined on this quantity.

\paragraph*{Statistical controls (Objection~10).}
Fold standard deviations alone understate variance, because the five folds
share one seed, one SAE, and one attack draw. Three controls calibrate what the
AUCs mean. A clean-vs-clean null, two independent realizations of the
preprocessing pipeline for the same images, gives AUC $0.500\pm0.018$ over $10$
repeats; a within-group label-permutation null gives $0.500\pm0.013$, placing
the observed $\varepsilon{=}0.004$ FGSM AUC of $0.932$ at $p=0.0099$ over $100$
permutations.
% src: experiments/2026-07-21_stat-package/outputs/stat_package.json (b, c blocks) [EST — E4b/E4c]
Headline $\dauc$ values also carry a paired bootstrap over source images
($95\%$ CI) and a five-seed replication with full SAE retraining and attack
regeneration (Section~\ref{sec:adv-specificity}).
% src: same file (a block); experiments/2026-07-02_gap-robustness/outputs/gap_robustness_summary.csv [EST — E4a/G5.3]
Two measurement bugs caught during this work motivated these controls; both are
documented as transferable case studies in Section~\ref{sec:adv-bugs}. For a
study whose central claims are that a control sits at chance and that a gap
vanishes, showing that the pipeline yields chance when it must is what makes
those nulls credible.

\paragraph*{Frozen-gate discipline.}
The pre-registered H1/H2 gates were fixed before any run and are never
re-evaluated on later data. Every experiment introduced after the gates failed
is exploratory: its expectations were written down before execution, and its
result is additive to the frozen verdict, never a substitute for it.
% src: EXPERIMENT_PLAN_ATTACKS.md (rev. 2026-07-21, rule block) [PRE/EXPL]

\subsection{Instrument scope: what pooling erases}
\label{sec:adv-instrument-scope}
Mean-pooling bounds the instrument, and the bound is measured rather than argued
about. A null result in this chapter is a null of a \emph{linear} probe on
\emph{mean-pooled} activations. Pooling over $197$ tokens erases any
token-localized perturbation by construction, and an earlier protocol showed
that a small non-linear probe on the same SAE latents recovered roughly $0.14$
AUC that the linear probe could not see. % src: D3b_RESULTS.md §3.1 [EST]
Consequently the chapter never claims that fine-grained adversarial structure is
\emph{absent}; it claims that such structure is invisible to a
magnitude-sensitive linear, mean-pooled probe under matched conditions. The
sharpest case, the steganographic control, is treated as Objection~8 in
Section~\ref{sec:adv-regime}, where the pooled null is rescoped to the
instrument by a token-level measurement.

% =============================================================================
\section{The pre-registered test: H1 and H2 at the frozen budget}
\label{sec:adv-gates}
% Objection 0. Both gates [PRE]; every number [EST].

The gates were defined as follows and never relaxed. \emph{H1 (magnitude
invariance)} passes if, over $\varepsilon\leq0.012$, the $\varepsilon$--AUC
correlation is near zero ($|r|<0.3$), the AUC variance is below $10^{-3}$, and
the mean AUC exceeds $0.95$. \emph{H2 (adversarial specificity)} passes if, at
the frozen budget $\varepsilon{=}0.012$ and against benign families of matched
nominal $L_2$, $\dauc>0.3$ and the strongest benign family is detected below
AUC $0.7$. % src: PLAN.md; DESIGN_CHOICES.md (frozen-gate definitions) [PRE]

\subsection{H1: detectability scales with the budget, for both attacks}
Table~\ref{tab:adv-h1} and Figure~\ref{fig:adv-h1} show detection AUC rising
monotonically with the perturbation budget for FGSM ($0.670$ to $1.000$,
saturating at $\varepsilon\geq0.012$) and for PGD, the pre-registered ``most
rigorous'' H1 test ($0.663$ to $1.000$, saturating earlier).
% src: both magnitude_sweep_metrics.csv [EST — L5]
This is the signature of a magnitude detector, and it fails the H1 gate on all
three criteria for both attacks: the correlation is strongly positive, the
variance is far above the threshold, and the mean over the five budgets
$\varepsilon\leq0.012$ is approximately $0.88$ for FGSM, below the floor.
% src: RESEARCH_LOG.md (corrected re-run verdict); PGD CSV [PRE — gate fails]
% The 0.88 is arithmetic on the five sourced FGSM values (0.670, 0.785, 0.932, 0.996, 1.000).

\begin{table}[ht]
  \centering
  \caption{H1: detection AUC and attack success rate (ASR) versus budget
  $\varepsilon$ under the corrected protocol, SAE latents at
  \texttt{mlp\_10}. Both attacks show the magnitude-detector signature: AUC
  rises monotonically and saturates ($\varepsilon\geq0.012$ for FGSM,
  $\approx0.008$ for PGD). The pre-registered H1 invariance gate fails all
  three criteria for both.}
  \label{tab:adv-h1}
  \begin{tabular}{lcccccccc}
    \toprule
    $\varepsilon$ & 0.001 & 0.002 & 0.004 & 0.008 & 0.012 & 0.016 & 0.024 & 0.032 \\
    \midrule
    FGSM AUC & 0.670 & 0.785 & 0.932 & 0.996 & 1.000 & 1.000 & 1.000 & 1.000 \\
    FGSM ASR & 0.382 & 0.619 & 0.762 & 0.819 & 0.832 & 0.831 & 0.817 & 0.808 \\
    PGD AUC  & 0.663 & 0.827 & 0.969 & 0.993 & 0.999 & 1.000 & 1.000 & 1.000 \\
    PGD ASR  & 0.580 & 0.922 & 0.999 & 1.000 & 1.000 & 1.000 & 1.000 & 1.000 \\
    \bottomrule
  \end{tabular}
  % src: experiments/2026-06-23_corrected-rerun/outputs/h1_sweep/magnitude_sweep_metrics.csv;
  %      experiments/2026-07-02_pgd-h1h2/outputs/h1_sweep/magnitude_sweep_metrics.csv
\end{table}

\begin{figure}[ht]
  \centering
  \includegraphics[width=0.8\linewidth]{figures/interp_adv/fig_h1_sweep.pdf}
  \caption{H1: detection AUC versus budget for FGSM and PGD, with
  success-conditioned AUC overlaid (Objection~2, Section~\ref{sec:adv-regime}):
  successful and failed attacks are equally detectable at every powered
  budget.}
  % src: the two magnitude_sweep_metrics.csv; experiments/2026-07-21_success-conditioned/outputs/metrics.json
  \label{fig:adv-h1}
\end{figure}

\subsection{H2: specificity at the frozen budget}
At the frozen budget the attack is at the ceiling and every benign family of
comparable nominal magnitude is detected at $0.85$--$0.94$ (JPEG $0.939$,
random dropout $0.921$, Gaussian blur $0.910$, Gaussian noise $0.847$, against
FGSM $0.9999$). The gap $\dauc = 0.061$ is measured against a ceiling and the
gate fails on both criteria: the $0.3$ dominance threshold is nowhere near
approached, and the strongest benign family far exceeds the $0.7$ bound.
% src: experiments/2026-06-23_corrected-rerun/outputs/h2_crossval/metrics.json [EST/PRE]
This run predates the per-family $L_2$ emission, so at the frozen budget the
$L_2$ match is the design intent written into the perturbation configurations;
the verified match arrives with the de-saturated re-run of
Section~\ref{sec:adv-regime}.

\subsection{Reading the failure in the right direction}
Both pre-registered gates fail, and fail stably across two attacks. That is the
intended negative outcome of the proposal, reported as such: $H_{M}$ is
eliminated in the tested regime. It would be a mistake, however, to read this
as a strong confirmation of $H_{D}$. Saturation \emph{removes} discriminative
power from the gates rather than adding it: once AUC is pinned at the ceiling,
the H1 gate can no longer separate a magnitude detector from a structure
detector, and the H2 gap is measured against a ceiling rather than in an
informative regime. Two rival readings therefore survive the frozen-budget
runs. The first is Objection~1: the gates fail only because the budget
saturates, and in a de-saturated regime the attack would separate from its
benign controls. The second is Objection~2: the monotone curve of
Table~\ref{tab:adv-h1} is equally compatible with ``the detector reads
adversarial structure, and structure scales with $\varepsilon$.'' Every
decisive test from here on runs in the informative sub-saturation regime
($\varepsilon\approx0.002$--$0.004$), and the next section takes up both
objections.

% =============================================================================
% Sections 4.5-4.11 written 2026-09-10 from the paper's final version (v06,
% 2026-07-28). Narrative contract: each section answers the objections listed
% in the ledger (Table~\ref{tab:adv-ledger}); refined and open entries are
% labelled as such in the prose. Notation follows the thesis convention: x is
% the input, a the probed (mean-pooled) activation, z / \hat{a} / r the SAE
% latents, reconstruction and residual; \Delta x is the activation displacement
% induced by perturbing x. Every number is copied from main.tex/appendix.tex;
% "% src:" trailers name the artifact.
% =============================================================================

\section{Objections about the regime and the reading}
\label{sec:adv-regime}
% Objections 1, 2 and 8. Objection 1 is NOT eliminated here: the de-saturated
% re-run produces a real residual gap, which is what forces the axis change of
% Section~\ref{sec:adv-specificity}.

The frozen-budget verdict of Section~\ref{sec:adv-gates} left two rival
readings standing, and the instrument scope of
Section~\ref{sec:adv-instrument-scope} raised a third. This section takes them
in turn. The first is answered by an experiment that made it stronger; the
chapter is organized around what that experiment forced.

\subsection{Objection 1: the gates fail only because the budget saturates}
\label{sec:adv-obj1}
At the frozen budget every representation detects FGSM at the ceiling, so the
H2 gap of $0.061$ is measured against a ceiling rather than in an informative
regime. The objection is that a de-saturated attack would separate from its
benign controls. The test re-runs H2 at $\varepsilon{=}0.004$, where FGSM is
detected at $0.932$ on the latents $z$ at \texttt{mlp\_10}, and where the
input-$L_2$ match is now \emph{verified} per family: the realized per-image
$L_2$ of every benign family lies within $0.999$--$1.008$ of the attack's
(Table~\ref{tab:adv-h2}).
% src: experiments/2026-06-29_h2-v2-eps0004/outputs/h2_crossval/metrics.json [EST]
The picture sharpens rather than collapses. The best benign family, Gaussian
blur, is detected at $0.733$; random dropout at $0.617$, JPEG at $0.604$, and
Gaussian noise at $0.574$. Hence $\dauc = 0.199$, with a paired bootstrap
$95\%$ interval of $[0.188, 0.212]$ over source images and a five-seed
replication, with full SAE retraining and attack regeneration, of
$0.204\pm0.014$. PGD at $\varepsilon{=}0.002$, its own sub-saturation budget,
gives the same: $\dauc = 0.224$, with every benign family below $0.61$.
% src: 2026-06-29 and 2026-07-02 h2_crossval/metrics.json; stat_package.json (a block); gap_robustness_summary.csv [EST]
The frozen gate still fails, since the $0.3$ dominance threshold is never
approached and the strongest benign family exceeds the $0.7$ bound; that
verdict is untouched. But a residual gap of $0.15$--$0.23$ is real by a
robustness criterion frozen after the first sub-saturated run and before the
five-seed replication that tests it ($\dauc - 2\sigma > 0.05$ at both
$\varepsilon\in\{0.002, 0.004\}$, at both mid and late layers, over five
seeds), and it has a depth profile: negative at \texttt{mlp\_0}, it rises to
$0.204\pm0.014$ at \texttt{mlp\_10} (Figure~\ref{fig:adv-decomp-layers}b).
% src: experiments/2026-07-02_gap-robustness/outputs/gap_robustness_summary.csv; DESIGN_CHOICES.md D12 (G5.3, frozen 2026-07-02, i.e. AFTER the 2026-06-29/30 eps=0.004 run and BEFORE the multiseed gap-robustness run) [PRE — robustness criterion; EST — values]

This is the strongest evidence against $H_{D}$ in the study, and the ledger
records it as such: Objection~1 is not eliminated by this experiment but
strengthened. An attack that is more detectable than every benign perturbation
of the same input norm looks like Degree-2 specificity. The reason it is not,
which is the hinge of the chapter, is that every control so far matches the
perturbation in \emph{input} space, while the probe reads \emph{activation}
space. If the attack and the benign families reach the probed layer with
different displacements, the residual gap has an entirely non-semantic
explanation. Section~\ref{sec:adv-specificity} measures the displacements and
shows that it does.

\begin{table}[ht]
  \centering
  \caption{H2 in the informative regime (SAE latents $z$, \texttt{mlp\_10},
  grouped cross-validation). The pre-registered gate ($\dauc>0.3$ and strongest
  benign family below $0.7$) fails here as at the frozen budget, yet a residual
  gap survives. ``$L_2$ ratio'' is the measured per-family input $L_2$ over the
  attack's; JPEG is absent from the PGD block because its realized $L_2$ cannot
  reach the $\pm20\%$ band. The bracketed interval is a paired bootstrap $95\%$
  confidence interval over source images.}
  % VERIFY (Breno): the ± values are carried from the paper's tab:h2 without a
  % stated meaning; confirm whether they are fold standard deviations before
  % describing them in the caption.
  \label{tab:adv-h2}
  \small
  \begin{tabular}{lcc}
    \toprule
    Family & AUC & $L_2$ ratio \\
    \midrule
    \multicolumn{3}{l}{\emph{FGSM, $\varepsilon{=}0.004$ (matched $L_2$ verified)}}\\
    FGSM                  & 0.932\,$\pm$\,0.005 & --- \\
    Gaussian blur         & 0.733\,$\pm$\,0.006 & 1.000 \\
    random dropout        & 0.617\,$\pm$\,0.005 & 1.005 \\
    JPEG                  & 0.604\,$\pm$\,0.006 & 1.008 \\
    Gaussian noise        & 0.574\,$\pm$\,0.007 & 0.999 \\
    $\dauc$               & \textbf{0.199} [0.188, 0.212] & \\
    \midrule
    \multicolumn{3}{l}{\emph{PGD, $\varepsilon{=}0.002$ (matched $L_2$ verified)}}\\
    PGD                   & 0.827\,$\pm$\,0.015 & --- \\
    Gaussian blur         & 0.603\,$\pm$\,0.007 & 1.000 \\
    random dropout        & 0.524\,$\pm$\,0.003 & 0.939 \\
    Gaussian noise        & 0.512\,$\pm$\,0.001 & 1.001 \\
    $\dauc$               & \textbf{0.224} & \\
    \bottomrule
  \end{tabular}
  % src: experiments/2026-06-29_h2-v2-eps0004/outputs/h2_crossval/metrics.json;
  %      experiments/2026-07-02_pgd-h2-eps0002/outputs/h2_crossval/metrics.json;
  %      experiments/2026-07-21_stat-package/outputs/stat_package.json
\end{table}

\subsection{Objection 2: adversarial structure grows with the budget}
\label{sec:adv-obj2}
The monotone curve of Table~\ref{tab:adv-h1} is compatible with a detector that
reads adversarial structure, provided that structure scales with
$\varepsilon$. If that were the mechanism, an attack that succeeds in flipping
the label, and therefore carries whatever structure does the flipping, should
be more detectable than one that fails at the same budget. Conditioning
detection on attack success removes the reading: successful and failed FGSM
examples are equally detectable at every powered budget, AUC $0.780$ against
$0.782$ at $\varepsilon{=}0.002$ ($N{=}950$ successful, $650$ failed), $0.928$
against $0.939$ at $0.004$ ($1170/430$), and $0.997$ against $0.999$ at
$0.008$ (the overlaid curves of Figure~\ref{fig:adv-h1}).
% src: experiments/2026-07-21_success-conditioned/outputs/metrics.json [EST — E3]
The detector responds to the perturbation's presence in activation space, not
to its functional effect. Objection~2 is eliminated. The result anticipates
the projection ablation of Section~\ref{sec:adv-residue}: what flips the label
and what fires the detector are largely orthogonal.

\subsection{Objection 8: the steganographic null shows a blind probe}
\label{sec:adv-obj8}
The steganographic control anchors the low-displacement column of the design
matrix, and the pooled probe cannot see it: AUC $0.542$ on $a$ and $0.519$ on
$z$ at \texttt{mlp\_10}, and on every representation it decays toward chance
with depth (Figure~\ref{fig:adv-decomp-layers}a). The objection is that this
null says nothing about displacement, because mean-pooling over $197$ tokens
erases a \emph{localized} perturbation by construction: a probe blind to the
watermark would sit at chance whatever the watermark did to the activation.
The objection is tested by probing token by token on the same activations
(Figure~\ref{fig:adv-stego-tokens}). The watermark's tokens light up: the best
single token reaches AUC $0.667\pm0.010$ (token~5; tokens $1$--$6$ span
$0.601$--$0.667$), against $0.542$ mean-pooled, $0.517$ on the class token
alone, and $0.512$ max-pooled.
% src: experiments/2026-07-21_stego-token/outputs/token_probe_metrics.json [EST — E9]
Activation space therefore does carry a trace of the watermark; pooling is what
hides it. And the pooled displacement is measured, not assumed: at
\texttt{mlp\_10} the control moves the representation by
$\lVert\Delta x\rVert=0.09$, under $1\%$ of the clean activation norm.
% src: experiments/2026-07-23_jpeg-matched-l2-dx/outputs/displacement_summary.json [EST]
This is exactly the behaviour a displacement-reading detector should show, a
local perturbation producing a small pooled displacement and no pooled signal,
so the null is rescoped to the instrument rather than discarded, and
Objection~8 is eliminated. The rescoping cuts both ways: no absence claim in
this chapter is a claim about token-resolved or non-linear probes
(Section~\ref{sec:adv-closing}).

\begin{figure}[ht]
  \centering
  \includegraphics[width=\linewidth]{figures/interp_adv/fig_decomp_layers.pdf}
  \caption{Detection across layers. (a) At the saturated budget
  ($\varepsilon{=}0.012$) every representation detects FGSM at the ceiling at
  every block, while the steganographic control decays to chance; nothing else
  can be read there. (b) In the informative regime ($\varepsilon{=}0.004$) the
  matched-input-$L_2$ gap $\dauc$ is a mid-to-late-layer phenomenon: negative or
  near zero at the first blocks, largest at \texttt{mlp\_10}. In the legend,
  $x$ denotes the raw activation $a$ of Equation~\eqref{eq:adv-decomp}.}
  % src: experiments/2026-06-22_sae-residual-decomp/outputs/decomposition_metrics.csv;
  %      experiments/2026-07-21_decomp-6layer/outputs/decomposition_metrics.csv;
  %      experiments/2026-07-02_gap-robustness/outputs/gap_robustness_summary.csv
  \label{fig:adv-decomp-layers}
\end{figure}

\begin{figure}[ht]
  \centering
  \includegraphics[width=0.7\linewidth]{figures/interp_adv/fig_stego_tokens.pdf}
  \caption{Per-token detection AUC for the steganographic control over the
  $14\times14$ patch grid at \texttt{mlp\_10}. The watermark's tokens (top
  left) reach $0.60$--$0.67$ while pooled probes sit near chance: the signal is
  local, and pooling, not activation space, is what hides it.}
  % src: experiments/2026-07-21_stego-token/outputs/token_probe_metrics.json
  \label{fig:adv-stego-tokens}
\end{figure}

% =============================================================================
\section{Objections about specificity: amplification and the matched-displacement control}
\label{sec:adv-specificity}
% Objections 3 and 4. This is the hinge of the chapter: the residual gap of
% Section~\ref{sec:adv-regime} is explained by amplification and vanishes once
% the benign families are matched on the axis the detector reads.

\subsection{Objection 3: the residual gap is Degree-2 specificity}
\label{sec:adv-obj3}
The residual $\dauc$ of $0.199$ survived every input-space control, and the
objection is that it is what it looks like: a signal specific to the
adversarial class. The test has two steps. The first measures, per example
and per layer, the displacement that each perturbation family induces at the
probed block, Equation~\eqref{eq:adv-dx}; the second re-runs H2 with the
benign families matched on that displacement rather than on the input.

\paragraph*{Adversarial optimization buys displacement cheaply.}
Matching input $L_2$ does not match what the detector reads
(Figure~\ref{fig:adv-amplification}). At matched input norm
($\varepsilon{=}0.004$), FGSM arrives at \texttt{mlp\_10} with a mean
displacement of $\lVert\Delta x\rVert=2.00$, while the four $L_2$-matched
benign families arrive at $0.19$--$0.59$ (Gaussian noise $0.19$, dropout
$0.21$, JPEG $0.28$, blur $0.59$); the supports barely overlap.
% src: experiments/2026-07-21_attack-lib-smoke/outputs/smoke_results.json; experiments/2026-07-23_jpeg-matched-l2-dx/outputs/displacement_summary.json [EST]
The clean activation norm at that layer is $10.0$, so the attack moves the
representation by $20\%$ of its length where benign perturbations of equal
input norm move it by $2$--$6\%$.
% src: displacement_summary.json (x_clean_norm_mean) [EST]
The Carlini--Wagner attack makes the same point from the other side. It was the
pre-registered candidate for the low-displacement adversarial cell of the design
matrix (Table~\ref{tab:adv-design}), because it minimizes input norm: it needs
only $L_2=4.97$, $0.72\times$ FGSM's budget at $\varepsilon{=}0.004$, to flip
every label. Yet it lands at $\lVert\Delta x\rVert=1.81\pm1.39$, attack-scale
displacement, and is detected at AUC $0.865$.
% src: experiments/2026-07-21_cw-attack/outputs/full/cw_input_l2.json; asr_cw.json; displacement_summary.json; eval/metrics.json [EST — E7]
PGD traces the same amplification along its sweep, from $\lVert\Delta x\rVert
= 1.24$ at $\varepsilon{=}0.001$ to $3.11$ at $0.032$; PGD at $0.002$ matches
FGSM at $0.004$ ($2.01$ against $2.00$), which is why those two budgets are
paired throughout.
% src: experiments/2026-07-21_pgd-displacement/outputs/pgd_displacement_summary.csv [EST — E6]
The low-displacement adversarial cell is therefore empty for every attack that
optimizes misclassification alone: even the input-norm-minimal attack is
amplified into a large displacement. The imperceptibility that defines these
attacks is an input-space property; in activation space they are loud, and that
is what makes them detectable. Adversarial optimization, precisely because it
pushes along directions the network is sensitive to, buys displacement cheaply
while benign corruptions do not. This is where adversarial structure enters
the account, as a \emph{cause} of displacement, upstream of anything the probe
reads.

\paragraph*{The matched-displacement control.}
The decisive experiment re-runs H2 with the four benign families calibrated,
model in the loop, to FGSM's displacement at the canonical layer (target
$1.994$; realized $1.78$--$2.06$; Section~\ref{sec:adv-attacks}). The result
is unambiguous (Table~\ref{tab:adv-matched-dx},
Figure~\ref{fig:adv-gap-regimes}): at matched $\lVert\Delta x\rVert$ the
FGSM--benign gap is \emph{negative or zero at every layer}, on the raw
activation and on the healthy-SAE latents alike, $\dauc^{z} = -0.062$ at
\texttt{mlp\_10} and $-0.078$ at \texttt{mlp\_11}, $\dauc^{a} = -0.011$ at
both, and $\dauc^{z}$ between $-0.000$ and $-0.108$ at the sub-gate early
layers.
% src: experiments/2026-07-22_h2-matched-dx/outputs/h2_matched_dx_summary.json [EST — E1]
A five-seed replication, with full SAE retraining and perturbation regeneration
per seed and the displacement match verified in all $25$ family--seed cells,
confirms it against a criterion fixed in advance: every cell stays negative,
including the $+2\sigma$ upper bounds ($\dauc^{z} = -0.063\pm0.003$ at
\texttt{mlp\_10}, $-0.076\pm0.005$ at \texttt{mlp\_11}). One flag is carried:
the dead-feature fraction exceeds $0.30$ at \texttt{mlp\_11} in two seeds, and
excluding those seeds leaves the verdict unchanged.
% src: experiments/2026-07-23_h2-matched-dx-multiseed/outputs/multiseed_summary.csv; multiseed_verdict.json; README.md (caveat) [EST — replicated 2026-07-23]

The input-$L_2$-matched gap of $0.199$ is thereby explained with no appeal to
adversarial meaning: it priced in the attack's privileged input-to-activation
amplification. Equalize the axis the detector reads and a blur is every bit as
``suspicious'' as the attack; the specificity that survived every input-space
control does not survive the displacement-space one. Objection~3 is
eliminated, and with it the Degree-2 reading. This is the hinge of the
chapter, and it is also the methodological contribution: the control that
standard detection setups omit is not a benign perturbation of matched input
norm, which the literature sometimes supplies
\cite{feinman2017adversarial_artifacts,ma2018lid_adversarial}, but a benign
perturbation matched on the displacement the detector reads.

\begin{table}[ht]
  \centering
  \caption{H2 with the benign families calibrated to FGSM's activation
  displacement (target $\lVert\Delta x\rVert=1.994$ at \texttt{mlp\_10};
  realized $1.78$--$2.06$). At matched displacement the best benign family is
  as detectable as, or more detectable than, the attack at every layer:
  $\dauc \leq 0$ throughout. $^{\dagger}$FVE below the $0.9$ health gate: $z$
  cells are reported for completeness only; $\dauc^{a}$ does not depend on the
  SAE. Five-seed replication: $\dauc^{z} = -0.063\pm0.003$ (\texttt{mlp\_10})
  and $-0.076\pm0.005$ (\texttt{mlp\_11}).}
  \label{tab:adv-matched-dx}
  \small
  \setlength{\tabcolsep}{4pt}
  \begin{tabular}{lccccc}
    \toprule
    Layer & FVE & FGSM $z$ & best benign $z$ & $\dauc^{z}$ & $\dauc^{a}$ \\
    \midrule
    \texttt{mlp\_0}  & 0.753$^{\dagger}$ & 0.944 & 0.999 & $-0.055$ & $-0.001$ \\
    \texttt{mlp\_3}  & 0.779$^{\dagger}$ & 1.000 & 1.000 & $-0.000$ & $0.000$ \\
    \texttt{mlp\_6}  & 0.720$^{\dagger}$ & 0.977 & 1.000 & $-0.023$ & $-0.007$ \\
    \texttt{mlp\_9}  & 0.697$^{\dagger}$ & 0.886 & 0.994 & $-0.108$ & $-0.032$ \\
    \texttt{mlp\_10} & 0.936 & 0.933 & 0.995 & $\mathbf{-0.062}$ & $\mathbf{-0.011}$ \\
    \texttt{mlp\_11} & 0.947 & 0.918 & 0.996 & $\mathbf{-0.078}$ & $\mathbf{-0.011}$ \\
    \bottomrule
  \end{tabular}
  % src: experiments/2026-07-22_h2-matched-dx/outputs/h2_matched_dx_summary.json;
  %      experiments/2026-07-23_h2-matched-dx-multiseed/outputs/multiseed_summary.csv
\end{table}

\begin{figure}[ht]
  \centering
  \includegraphics[width=0.85\linewidth]{figures/interp_adv/fig_gap_regimes.pdf}
  \caption{$\dauc$ across matching regimes (\texttt{mlp\_10}, latents $z$).
  Measured against a ceiling at the frozen budget, real under verified matched
  input-$L_2$, it inverts once the benign families are matched on the axis the
  detector reads. Error bars: paired bootstrap $95\%$ confidence interval
  (matched $L_2$); $\pm2\sigma$ over five seeds (matched
  $\lVert\Delta x\rVert$).}
  % src: gate blocks of the three h2 metrics.json; multiseed_summary.csv; stat_package.json
  \label{fig:adv-gap-regimes}
\end{figure}

\subsection{Objection 4: the SAE is load-bearing}
\label{sec:adv-obj4}
The objection is that the gap lives in the interpretable latents $z$ rather
than in the raw activation $a$, so that the SAE, not displacement, is what
exposes the attack. Read after the matched-displacement verdict, the residual
decomposition of Equation~\eqref{eq:adv-decomp} is a post-mortem on the
input-matched gap, and it settles two things.

First, detection needs no SAE. At $\varepsilon{=}0.004$ the raw activation $a$
detects FGSM at $0.989$ against $0.932$ on the latents $z$
(Table~\ref{tab:adv-decomp004}), and in the saturated regime every
representation is at the ceiling at every layer
(Table~\ref{tab:adv-decomp}). In information terms the SAE is a lossy
bottleneck, so $z$ cannot carry detection information that $a$ lacks; a linear
probe on $z$ could still profit from the encoder's non-linearity, but here it
does not, and for detection the SAE is not load-bearing.
% src: experiments/2026-06-30_decomp-eps0004/outputs/decomposition_metrics.csv; experiments/2026-06-22_sae-residual-decomp/outputs/decomposition_metrics.csv [EST — L2]
Second, the input-$L_2$-matched gap concentrates in the SAE-explained,
on-manifold component ($0.199$ on $z$, $0.178$ on $\hat{a}$) rather than in the
off-manifold residual ($0.079$ on $r$, $0.066$ on $a$). The generic detection
signal lives in $a$ and $r$, which detect at $0.989$--$0.990$ with small gaps,
while the apparent family separation lives on the manifold, exactly where the
amplification premium lands when the wrong axis is matched. The scalar residual
norm $\lVert r\rVert$ collapses to $\dauc = 0.054$ even as the residual
\emph{vector} stays at $0.990$: the probe reads displacement direction, not
size. An eleven-signal sweep over the same cell (Figure~\ref{fig:adv-zoo})
confirms the ordering, with per-feature SAE codes leading the input-matched gap
and scalar magnitudes trailing.
% src: 2026-06-30_decomp-eps0004 decomposition_metrics.csv; 2026-07-02_signal-zoo/outputs/zoo/zoo_metrics.csv [EST]

Two qualifications keep this honest. The larger gap on $z$ under input
matching is real, and a pre-registered anti-condition forbids over-reading it:
at $\varepsilon{=}0.004$ the logit-scale difference on the raw activation $a$
is compatible with that on $z$, so ``the SAE amplifies specificity'' may not be
claimed from that cell (only at $\varepsilon{=}0.002$ is $z$ logit-distinct
from $a$).
% src: gap_robustness_summary.csv (logit columns); DESIGN_CHOICES.md D12 (G5.3, frozen 2026-07-02) [PRE — anti-condition]
And the gap on $z$ has the depth profile of Section~\ref{sec:adv-obj1}, absent
at \texttt{mlp\_0} and largest at \texttt{mlp\_10}, whereas displacement
detectability is present already at the first block. Once the displacement is
equalized, none of the on-manifold separation survives
(Table~\ref{tab:adv-matched-dx}). Objection~4 is eliminated for detection: the
SAE's interpretable features are a lossy view of a signal that the raw
activation carries better, and the one thing they appeared to add was
amplification read on the wrong axis.

\begin{table}[ht]
  \centering
  \caption{Residual decomposition at the sub-saturated budget
  ($\varepsilon{=}0.004$, \texttt{mlp\_10}, $N{=}3200$, SAE FVE $0.938$).
  $\dauc$ (Equation~\eqref{eq:adv-dauc}) is the AUC on FGSM minus that on the
  strongest \emph{matched input-$L_2$} benign family, Gaussian blur for every
  representation. The generic detection signal lives in $a$ and $r$ (high AUC,
  small gap); the family separation under input-$L_2$ matching concentrates in
  the SAE-explained component $z$, $\hat{a}$, and vanishes at matched
  displacement (Table~\ref{tab:adv-matched-dx}).}
  \label{tab:adv-decomp004}
  \begin{tabular}{lcccc}
    \toprule
    Representation & FGSM & Blur (best benign) & Stego & $\dauc$ \\
    \midrule
    $a$              & 0.989 & 0.923 & 0.542 & 0.066 \\
    $z$              & 0.932 & 0.733 & 0.519 & 0.199 \\
    $\hat{a}$        & 0.963 & 0.785 & 0.532 & 0.178 \\
    $r$              & 0.990 & 0.911 & 0.539 & 0.079 \\
    $\lVert r\rVert$ & 0.594 & 0.540 & 0.502 & 0.054 \\
    \bottomrule
  \end{tabular}
  % src: experiments/2026-06-30_decomp-eps0004/outputs/decomposition_metrics.csv
\end{table}

\begin{table}[ht]
  \centering
  \caption{Detection AUC (grouped cross-validation, $K{=}5$, $N{=}3200$) by
  representation and layer at the \emph{saturated} budget
  ($\varepsilon{=}0.012$). FGSM is at the ceiling everywhere; the
  steganographic control is near chance. At the ceiling cross-representation
  comparisons are vacuous, which is what motivates the sub-saturated regime of
  Table~\ref{tab:adv-decomp004}.}
  \label{tab:adv-decomp}
  \small
  \begin{tabular}{llcccc}
    \toprule
    Layer & Task & $a$ & $z$ & $\hat{a}$ & $r$ \\
    \midrule
    \multirow{2}{*}{\texttt{mlp\_0}}  & FGSM  & 1.000 & 0.996 & 1.000 & 1.000 \\
                                      & stego & 0.675 & 0.539 & 0.541 & 0.677 \\
    \multirow{2}{*}{\texttt{mlp\_3}}  & FGSM  & 1.000 & 1.000 & 1.000 & 1.000 \\
                                      & stego & 0.685 & 0.604 & 0.604 & 0.664 \\
    \multirow{2}{*}{\texttt{mlp\_6}}  & FGSM  & 1.000 & 1.000 & 1.000 & 1.000 \\
                                      & stego & 0.549 & 0.531 & 0.537 & 0.542 \\
    \multirow{2}{*}{\texttt{mlp\_9}}  & FGSM  & 1.000 & 0.999 & 0.998 & 1.000 \\
                                      & stego & 0.530 & 0.513 & 0.524 & 0.530 \\
    \multirow{2}{*}{\texttt{mlp\_10}} & FGSM  & 1.000 & 1.000 & 1.000 & 1.000 \\
                                      & stego & 0.543 & 0.519 & 0.532 & 0.540 \\
    \multirow{2}{*}{\texttt{mlp\_11}} & FGSM  & 1.000 & 0.997 & 1.000 & 1.000 \\
                                      & stego & 0.544 & 0.513 & 0.524 & 0.542 \\
    \bottomrule
  \end{tabular}
  % src: experiments/2026-06-22_sae-residual-decomp/outputs/decomposition_metrics.csv
\end{table}

\begin{figure}[ht]
  \centering
  \includegraphics[width=0.85\linewidth]{figures/interp_adv/fig_zoo.pdf}
  \caption{The signal zoo: input-$L_2$-matched $\dauc$ (FGSM minus the
  strongest matched benign family) for eleven detection signals at
  $\varepsilon{=}0.004$, \texttt{mlp\_10}. Per-feature SAE codes lead; scalar
  magnitudes trail. In the legend, $x$ and $\hat{x}$ denote the raw activation
  $a$ and its reconstruction $\hat{a}$ of Equation~\eqref{eq:adv-decomp}.}
  % src: experiments/2026-07-02_signal-zoo/outputs/zoo/zoo_metrics.csv
  \label{fig:adv-zoo}
\end{figure}

% =============================================================================
\section{What remains at equal displacement}
\label{sec:adv-residue}
% Objections 5, 6 and 9. Objection 5 is REFINED, not eliminated: the families
% do not collapse onto one curve, but the axis that separates them is texture.
% Objection 9 is REFINED: Degree 3 is a learned direction, not any distance.

With the Degree-2 reading gone, the remaining objections concern what the
detector reads at equal displacement. They are answered by three experiments,
two of which do not end in elimination but in a refinement of $H_{D}$.

\subsection{Objection 5: if only magnitude mattered, all families would collapse onto one curve}
\label{sec:adv-obj5}
The strongest form of $H_{D}$ predicts \emph{collapse}: a single detector's
per-example score should be a function of $\lVert\Delta x\rVert$ alone, one
curve for all families. The test is direct. One clean-versus-FGSM logistic
detector on $z$ at \texttt{mlp\_10} (held-out AUC $0.933$) scores the four
benign families out of sample; per-example scores are paired to per-example
displacements; and a nested $F$-test asks whether family identity adds
explanatory power over a polynomial in $\lVert\Delta x\rVert$.
% src: experiments/2026-07-22_collapse-e5/outputs/metrics.json (protocol block) [EXPL — E5, expectations pre-stated]
Collapse \emph{fails}: family adds power ($\Delta R^2=0.373$, $F=1204$,
$N=8000$, $p\approx0$).
% src: same file (collapse_test) [EST]
But the way it fails is the finding (Figure~\ref{fig:adv-collapse}). In the
matched band $\lVert\Delta x\rVert\in[1.8,2.2]$ the mean scores are FGSM
$0.877$, Gaussian noise $0.890$, random dropout $0.519$, JPEG $0.291$, and
Gaussian blur $0.206$.
% src: same file (per_family_score, matched band) [EST]
The attack's nearest neighbour on the score axis is a \emph{benign} family,
i.i.d.\ noise, while the smooth corruptions score low at the same displacement.
The axis that separates families is not adversarial-versus-benign but
high-frequency-noise-like versus smooth: \emph{texture}. The detector is not
family-blind, but the information it uses is textural, and the attack carries
no premium over benign noise of equal displacement and similar texture.

Objection~5 therefore refines $H_{D}$ rather than eliminating it, and the
ledger records the amendment. Together with the matched-displacement verdict,
the refined claim reads: the detector reads off-manifold displacement magnitude
plus perturbation texture; adversarial meaning enters only upstream, as the
cause of the displacement, not as something the probe reads.

\begin{figure}[ht]
  \centering
  \includegraphics[width=0.9\linewidth]{figures/interp_adv/fig_collapse.pdf}
  \caption{The collapse test. Per-example score of a single clean-versus-FGSM
  detector against $\lVert\Delta x\rVert$ at \texttt{mlp\_10}; dots are
  examples, lines are binned means. The curves do not collapse, but the axis
  that separates them is texture: FGSM rides the same curve as benign Gaussian
  noise, far above the smooth and structured families at equal displacement.}
  % src: experiments/2026-07-22_collapse-e5/outputs/collapse_table.csv
  \label{fig:adv-collapse}
\end{figure}

\subsection{Objection 6: the probe reads the direction that flips the label}
\label{sec:adv-obj6}
If the probe read the \emph{attack}, it would read the part of $\Delta x$
inside the classification head's sensitive subspace. At the representation the
linear head consumes this subspace is exact, $J = W$, so the displacement
splits as $\Delta x = \Delta x_\parallel + \Delta x_\perp$ into the functional
direction and its orthogonal complement.
% src: EXPERIMENT_PLAN_ATTACKS.md §E12 (definition) [EXPL — E12a]
On FGSM at $\varepsilon{=}0.004$ ($N{=}1600$), $87.7\%$ of the displacement's
energy is in $\Delta x_\perp$, and a projection ablation shows that the
detector lives there: AUC $0.968$ on the full displacement, $0.967$ on
$\Delta x_\perp$ alone, and $0.694$ on $\Delta x_\parallel$ alone
(Figure~\ref{fig:adv-jspace}).
% src: experiments/2026-07-21_jacobian/outputs/e12a_full/jacobian_metrics.json [EST]
Two scoping notes keep this honest. The head has rank at most $10$ in $768$
dimensions, so an isotropic displacement would put only about $1.3\%$ of its
energy in $\Delta x_\parallel$; the measured $12.3\%$ is a roughly $10\times$
\emph{enrichment}. The attack does concentrate energy in the functional
direction, and the probe reads mostly $\Delta x_\perp$ \emph{despite} that, not
because $\Delta x_\parallel$ is empty. And the ablated AUCs compare subspaces of
very different dimension (at most $10$ against about $758$), so they locate
where the signal lives rather than forming a paired contrast.
% [EST — arithmetic on sourced values: 10/768 = 0.013; 1 - 0.877 = 0.123]
Objection~6 is eliminated with that caveat. What flips the label and what
fires the detector are largely orthogonal, consistent with detection being
independent of attack success (Section~\ref{sec:adv-obj2}) and with
amplification: the attack pushes through the head via a small
$\Delta x_\parallel$, while its large, detectable $\Delta x_\perp$ wake is what
the probe sees.

\begin{figure}[ht]
  \centering
  \includegraphics[width=0.6\linewidth]{figures/interp_adv/fig_jspace.pdf}
  \caption{Projection ablation at the head input ($J = W$, exact). Removing the
  functional direction $\Delta x_\parallel$ costs the detector nothing; the
  functional direction alone is barely detectable.}
  % src: experiments/2026-07-21_jacobian/outputs/e12a_full/jacobian_metrics.json
  \label{fig:adv-jspace}
\end{figure}

\subsection{Objection 9: then any off-manifold distance should reproduce the probe}
\label{sec:adv-obj9}
Degree~3 does not mean that any off-manifold distance fires. A Mahalanobis
distance in the spirit of \citeonline{lee2018mahalanobis}, computed on the same
mean-pooled raw activations under the same source-image splits, does
\emph{not} reproduce the trained probe: on FGSM at $\varepsilon{=}0.004$,
\texttt{mlp\_10}, it reaches $0.710$ where the logistic probe reaches $0.989$,
and it is at or near chance on the matched benign families ($0.482$--$0.557$)
and on the steganographic control ($0.498$).
% src: experiments/2026-07-21_mahalanobis/outputs/maha_vs_logit.csv; maha_baseline.py (single clean Gaussian, shrinkage 1e-3, per-layer) [EST — E8]
The baseline is run in its simplest form, a single Gaussian fitted to the clean
training-fold activations at one layer, without the original's layer-wise
ensemble or input pre-processing, so the claim is not that the method of
\citeonline{lee2018mahalanobis} fails but that a bare off-manifold distance does
not land where the probe does. Objection~9 refines $H_{D}$ a second time: what
the chapter characterizes is a \emph{learned} direction in displacement space,
one that aligns with displacement magnitude and texture rather than with
adversarial identity. The detector is generic in \emph{what} it responds to
and specific in \emph{how} it reads it.

% =============================================================================
\section{The empty cell: a minimal-displacement attack}
\label{sec:adv-lowdx}
% Objection 7 — OPEN. Exploratory run (E7 lineage, 2026-07-25), outcome
% branches and benign band pre-stated before the full run. This is NOT the
% frozen A5.4 adaptive PGD, which remains unrun (Section~\ref{sec:adv-closing}).
% Decision (Breno, 2026-07-28): the asymmetry is reported, not interpreted.

The last objection is the deepest: adversarial and high-displacement are
confounded by construction, because every attack examined so far optimizes
misclassification and pays displacement as collateral. If the top-right cell
of the design matrix (Table~\ref{tab:adv-design}) cannot be occupied, $H_{D}$
has never been tested against a successful attack at benign displacement. The
section attempts to occupy the cell on purpose.

``Quiet'' means two incompatible things here. An adversarial example is quiet
\emph{to the human}: its input perturbation is small and the image is visually
unchanged. It is anything but quiet \emph{to the detector}: the same
imperceptible perturbation moves the activation by $2.00$, and that is why the
detector works. The empty cell asks for the opposite trade, an attack quiet to
the detector whatever it costs the input. The attack is run with that
objective, per image minimizing the displacement at the canonical block,
\begin{equation}
  \min_{\tilde{x}} \;\lVert a_{10}(\tilde{x}) - a_{10}(x)\rVert_2^2
  \;+\; c\cdot\max\bigl(\mathrm{margin}(\tilde{x}), -\kappa\bigr),
  \label{eq:adv-mindisp}
\end{equation}
where $a_{10}$ is the mean-pooled \texttt{mlp\_10} activation, a $\tanh$
reparametrization keeps $\tilde{x}$ in pixel space, and an outer binary search
drives $c$ to the smallest weight that still flips the label. The \emph{floor}
is the smallest realized $\lVert\Delta x\rVert$ among successful iterates, an
upper bound on the true minimum by construction. The benign band was fixed in
advance at $\lVert\Delta x\rVert \leq 0.30$, anchored on the measured
displacement of matched-$L_2$ JPEG, $0.284$.
% src: experiments/2026-07-25_min-displacement-cw/README.md (pre-registration block); src/adv_interp/attacks/cw_min_displacement.py [EXPL — outcome branches pre-stated]

\paragraph*{The cell is occupiable, at a price.}
On all $1536$ images that the model classifies correctly when clean, the attack
succeeds (ASR $1.000$) with a floor of median $\lVert\Delta x\rVert = 0.298$ at
\texttt{mlp\_10} (interquartile range $[0.235, 0.368]$; $50.7\%$ of the images
inside the benign band), an order of magnitude below FGSM at
$\varepsilon{=}0.004$ ($2.00$) and six times below input-norm-minimal CW
($1.81$); one block up, at \texttt{mlp\_11}, the same images sit at $0.579$.
% src: experiments/2026-07-25_min-displacement-cw/outputs/full/floor_summary.json; verdict.json [EST — pre-registered branch 1, marginally: the floor straddles the benign anchor]
A successful attack at benign displacement scale exists; finding it required
optimizing displacement itself. The price is the other kind of silence, and it
is not subtle: holding the activation nearly still while flipping the label
costs a median pixel $L_\infty$ of $0.137$, $34\times$ FGSM's budget at
$\varepsilon{=}0.004$ ($L_2 = 8.39$), on visibly perturbed images.
% src: same floor_summary.json (input_linf/input_l2 blocks) [EST]
This is the section's result, not a caveat on it: the two silences are
mutually exclusive here. Input norm and activation displacement trade off, so
quiet to the detector is bought with loud to the eye, exactly as quiet to the
eye is bought with loud to the detector; an attacker gets one form of
imperceptibility or the other, not both. The floor is an upper bound, and a
fair one: the search bottoms out rather than exhausting its budget (median
successful displacement flat at about $0.31$ over the last three halvings of
$c$), though at roughly half the sketched inner budget, so the safe statement
is that no success was found below the first quartile, about $0.24$.
% src: same experiment, outputs/full/pareto_trace.csv; params.yaml (attack.search_budget_rationale) [EST — upper-bound discipline]

\paragraph*{What the detector does on the floor set.}
The detector is \emph{not} evaded: clean-versus-floor AUC is $0.999$ on $a$
and $0.990$ on $z$ at \texttt{mlp\_10}.
% src: outputs/full/detection_auc.json (floor_set cells) [EST]
This is the run's sharpest fact, because at essentially the same displacement
the benign control sits far lower: JPEG at $\lVert\Delta x\rVert = 0.284$ is
detected at $0.604$ on $z$, an asymmetry of about $0.39$ AUC at matched
displacement.
% src: detection_auc.json; 2026-06-29_h2-v2-eps0004/outputs/h2_crossval/metrics.json (benign_jpeg) [EST — the asymmetry itself]
The asymmetry is reported without being interpreted, because one alternative
explanation stays open: the floor set's input perturbation is enormous
compared with JPEG's, so its small displacement may simply point far off the
data manifold, a \emph{direction} that a linear probe reads without reading
magnitude. Separating ``adversarial structure beyond displacement'' from
``input-driven off-manifold direction'' needs a control this run lacks, a
benign family matched to the floor set's input $L_\infty$. Objection~7 is
therefore left as a bounded open question in the ledger: it marks the edge of
the displacement reading rather than overturning it, and the missing control
is the natural next experiment.
% [DECISION 2026-07-28 (Breno): reported, not interpreted; L∞-matched control not run]
The mechanism reading, meanwhile, is unchanged and sharper than for FGSM:
$92.4\%$ of the floor displacement's energy is orthogonal to the head's
functional subspace (FGSM: $87.7\%$), the probe's signal lives there (AUC
$0.9996$ on the full displacement, $0.9997$ on $\Delta x_\perp$ alone, $0.7804$
on $\Delta x_\parallel$ alone), and locally at \texttt{mlp\_10} the functional
component alone is at chance ($0.518$).
% src: outputs/full/perp_parallel.json [EST — E12-consistent]
Even an attack optimized to move the activation as little as possible is read
by the part of its displacement that is not the attack.

% =============================================================================
\section{Reconciliation with deployment work}
\label{sec:adv-reconciliation}
% Objection 11 — SAEgis. All three corroborations verified against
% arXiv:2605.07447 (HTML) on 2026-07-24; paraphrase only.

The remaining objection is external: SAEgis \cite{wang2026saegis} builds a
working SAE-based adversarial firewall for vision--language models, so, the
objection runs, SAE-based detection works. The specificity ladder of
Section~\ref{sec:adv-ladder} locates the disagreement. SAEgis's own
cross-attack analysis already argues against Degree~1: the features its
firewall uses overlap heavily across attack methods at the earliest vision
block, with the overlap decreasing with depth.
% [EXT — verified: near-complete overlap between attack families at vision block 0; claim scoped to that block]
The matched input-$L_2$ control of this chapter is the Degree-2 instrument:
taken alone it finds \emph{apparent} class-specificity
(Section~\ref{sec:adv-regime}), and the displacement axis dissolves it
(Section~\ref{sec:adv-specificity}). The detector sits at Degree~3, the degree
an input-space analysis alone could never reach, and Section~\ref{sec:adv-obj9}
qualifies what Degree~3 means.

This makes the chapter's result a complement to deployment work, not a
contradiction of it: none of SAEgis's measurements is contradicted, and where
the two studies overlap they corroborate this one. Its dense, SAE-free
baseline ties the SAE firewall in domain and beats it on one domain, matching
$a \geq z$ here. Its SAE reconstruction-error score, which is the scalar
$\lVert r\rVert$ of this chapter up to a monotone transform (per-example
reconstruction error equals $\lVert r\rVert^2/d_{\mathrm{model}}$), gives a
negligible clean--adversarial difference, matching $\dauc = 0.054$ on
$\lVert r\rVert$ (Table~\ref{tab:adv-decomp004}). And its earliest vision
block, its most robust cross-attack detector, is the high-frequency-texture
signal that the matched-displacement analysis isolates
(Section~\ref{sec:adv-obj5}).
% [EXT — all three verified 2026-07-24: (1) dense baseline = cosine to averaged embeddings on raw hidden states, ties on two domains and beats SAEgis on the medical domain; (2) reconstruction-MSE anomaly score shows negligible clean-vs-adversarial difference; (3) vision block 0 most robust cross-attack, high-frequency-texture hypothesis]
What is contested is only the strong reading, that such detectors read
adversarial-class-specific meaning; the weak reading, that the features build a
useful detector, stands. Objection~11 is repositioned rather than eliminated:
this chapter is, in effect, the control experiment that deployment papers
need, at both degrees, the input-norm match and the displacement match that
shows even the former is not control enough.

% =============================================================================
\section{Methodological case studies}
\label{sec:adv-bugs}
% Two measurement bugs, kept in the chapter because they motivated the
% statistical controls of Section~\ref{sec:adv-protocol}. May move to an
% appendix at final assembly.

Two measurement bugs caught during this work are documented here because each
produced a plausible result that was wrong, and each yields a transferable
rule.

\paragraph*{The clip bug: an out-of-distribution artefact masquerading as
magnitude-invariant detection.}
The FGSM routine originally applied $\mathrm{clamp}(\cdot, 0, 1)$ to
ImageNet-\emph{normalized} images, whose values span roughly $[-2.2, 2.6]$.
% src: DESIGN_CHOICES.md; RESEARCH_LOG.md 2026-04-23 [EST]
Clipping in that space squashes most pixels: at $\varepsilon{=}0.012$ the
measured per-image $L_2$ was about $284$ against the $20$ that the
$\varepsilon$ arithmetic implies, and the clipped images left the normalized
manifold entirely, a synthetic, $\varepsilon$-independent out-of-distribution
signal that the detector simply read off. The pre-fix sweep was therefore
\emph{flat} at AUC $\approx 0.88$, a false confirmation of H1: the artefact
contributed about $0.23$ AUC at $\varepsilon{=}0.001$ but only $0.02$ at
$0.032$, cancelling the true magnitude trend.
% src: RESEARCH_LOG.md 2026-04-23 (delta table) [EST]
The fix clips where it is physically meaningful (denormalize, clip in $[0,1]$,
renormalize); post-fix $L_2$ at $\varepsilon{=}0.012$ is $20.51$, consistent
with theory, and the sweep recovers the monotone curve of
Table~\ref{tab:adv-h1}.
% src: src/adv_interp/attacks/fgsm.py [EST]
The transferable rule: measure the realized perturbation norm against the
budget arithmetic \emph{before} reporting any detection number. Had the bug
survived, the flat sweep would have satisfied the two invariance criteria of
the pre-registered H1 gate (correlation and variance) on an artefact, failing
only the mean-AUC floor.

\paragraph*{The first-key bug: a layer-provenance error in SAE training.}
The SAE trainer selected its target activations as the first non-label key of
the saved activation file. With forward hooks registered in the order
$[\texttt{layer\_10}, \texttt{layer\_11}]$, that first key was always
\texttt{layer\_10}, while the \texttt{target\_layer} field in the configuration
files was dead code.
% src: CLUSTER_NOTES.md 2026-06-22; RESEARCH_LOG.md [EST]
The project's ``canonical'' detection number was thus computed on
\texttt{mlp\_10} alone, not on the \texttt{mlp\_11} or combined layers the
documentation claimed: a provenance error, not an invalidation of the
underlying measurement, since the layer-sweep experiment, which isolates one
layer per run, was unaffected. The fix plumbs an explicit
\texttt{--target-layer} argument through SAE training and latent encoding,
removing the silent dependence on hook registration order; \texttt{mlp\_10} was
then adopted as the canonical layer on the record.
% src: CLUSTER_NOTES.md (P5 rewire) [EST]
The transferable rule: never let an analysis target be selected implicitly by
container ordering; make it an explicit, logged argument.

% =============================================================================
\section{Limitations and the reading for the thesis}
\label{sec:adv-closing}

\subsection{Limitations}
The claims of this chapter are bounded by the instrument and by the scope of
the runs behind them.

\paragraph*{A null is a null of this probe.}
Detection here is a \emph{linear} classifier on \emph{mean-pooled}
activations, so every absence claim is scoped to that instrument. Pooling can
hide localized signal (Section~\ref{sec:adv-obj8}), and a small non-linear
probe once recovered about $0.14$ AUC that the linear probe could not see under
an earlier protocol (Section~\ref{sec:adv-instrument-scope}). A token-level,
non-linear probe remains future work.
% src: token_probe_metrics.json; zoo_metrics.csv; D3b_RESULTS.md §3.1 [EST]

\paragraph*{Residual single-seed scope.}
The matched-displacement control, the hinge of the chapter, is replicated
across five seeds; its displacement calibration is performed at
\texttt{mlp\_10} and inherited by the other layers. The collapse test, the CW
attack, the minimal-displacement attack, the token-resolved probe, and the
projection ablation remain single-seed (the last two single-layer), and the
six-layer decomposition ran at reduced sample.
% src: 2026-07-23 multiseed; 2026-07-22_collapse-e5; 2026-07-21_cw-attack; 2026-07-25_min-displacement-cw; 2026-07-21_stego-token; 2026-07-21_jacobian; 2026-07-21_decomp-6layer [EST — scope]

\paragraph*{SAE health limits latent claims.}
The FVE $\geq 0.9$ gate passes only at \texttt{mlp\_10} and \texttt{mlp\_11},
so all claims about $z$ are scoped to those layers; the per-layer FVE also
disagrees between runs at early layers, an unexplained fit difference that is
flagged rather than resolved (Section~\ref{sec:adv-sae}), and the dead-feature
fraction exceeds its bound at \texttt{mlp\_11} for two seeds, whose exclusion
leaves the verdict unchanged. Healthy per-layer SAEs are future work.
% src: h2_matched_dx_summary.json; 2026-07-21_decomp-6layer; 2026-07-23 multiseed README [EST]

\paragraph*{Coverage.}
One dataset (CIFAR-10 upsampled over an ImageNet-pretrained backbone, with no
native-resolution replica) and four attack constructions plus the
steganographic control. The minimal-displacement attack is single-seed and its
asymmetry awaits the input-$L_\infty$-matched control; the pre-registered
\emph{adaptive} PGD (item A5.4 of the plan) remains unrun, so nothing here
speaks to an adversary that optimizes against the detector; and the ViT-L/16
scale check ran only under an earlier protocol. The system is a ViT
classifier, and no claim is made about vision--language models.
% src: DESIGN_CHOICES.md D12 (A5.4); 2026-07-25_min-displacement-cw; 2026-07-23_jpeg-matched-l2-dx [PRE/EXPL/scope]

\subsection{The reading for the thesis}
The chapter set out to characterize what an activation-space adversarial
detector in a ViT classifier measures. The pre-registered hypotheses fail, and
fail stably: detection AUC scales monotonically with budget for FGSM and PGD
alike, and no regime tested, saturated or de-saturated, either attack, brings
the adversarial--benign gap near its frozen threshold. The contribution beyond
the negative is what the residual turned out to be. The $\dauc$ of $0.199$
that survived every input-space control does not survive the displacement one:
measured on the axis the detector reads, the apparent specificity is
amplification, not meaning. This is the sense in which detection here is
\emph{magnitude, not meaning}: magnitude means measured activation
displacement, and the attack's only privileged role is how cheaply its
optimization buys it. The two axes trade off, since a classical adversarial
example is imperceptible in input space and, precisely because it is, loud in
activation space, while a purpose-built attack reaches a benign corruption's
displacement only at a $34\times$ input-norm premium on visibly corrupted
images. The asymmetry that attack leaves behind is owned as the edge of the
reading, not folded into it.

For the thesis, three things follow. First, the diagnostics question of
Chapter~\ref{chpt:intro} receives a definite answer in this setting: for
perturbations that optimize misclassification alone, sparse-feature signatures
do not isolate a structurally specific perturbation; they collapse into change
detection, and the change is measurable as one physical variable per example.
Second, the sparse autoencoder is not load-bearing for that detection. Its
latents are a lossy view of a signal that the raw activation carries better,
and the one thing they appeared to add, a larger gap under input-norm
matching, was amplification read on the wrong axis. The interpretable
substrate is a legitimate measuring instrument, since it defines the
on-/off-manifold split, but it is not the source of the signal. Third, and
this is the part that travels beyond this system, the result was only visible
because of a control that standard evaluations omit and that the
pre-registration itself did not anticipate: benign perturbations matched on
the displacement the detector reads, not on the input. The frozen gates were
necessary to keep the negative honest, but they were not sufficient to find
the explanation; that took an exploratory experiment whose expectations were
written down before it ran and whose status the ledger keeps distinct from
the pre-registered verdict. Before reading meaning into an activation-space
detector, measure the per-example displacement it induces and match the
controls there. Matching the input is not control enough.

\end{document}
